% GENERATED FILE. Edit canonical lessons, then run npm run generate:books.
% canonical-source-hash: fnv1a64:4ae5b04642f0c3e1
% canonical-lessons: ML-C273-perukkuka, ML-C273-vitaykkuka, ML-C273-kottuka, ML-C273-karannuka, ML-C273-nulluka

\chapter{To pick up, To sow, To peck, To spin, To pinch}
\label{ch:ml-to-pick-up-to-sow-to-peck-to-spin-to-pinch}
\hlchaptermodality{273}

\begin{chapteropening}
\textbf{By the end of this chapter:} \emph{I can say to pick up, to sow, to peck, to spin and to pinch.}

Complete the last lesson of chapter 273: I can say to pick up, to sow, to peck, to spin and to pinch.
\end{chapteropening}

\section[\texorpdfstring{\ml{പെറുക്കുക} (peṟukkuka)}{പെറുക്കുക (peṟukkuka)}]{\ml{പെറുക്കുക} (peṟukkuka) — to pick up (from the ground)}

\label{lesson:ML-C273-perukkuka}



\begin{quote}
Before the new one: say the Malayalam for to drill, then the Malayalam for to leak.
\end{quote}

\subsection*{You'll want to know: \ml{പെറുക്കുക}}
\textbf{\ml{പെറുക്കുക}} — \emph{peṟukkuka} — \textquotedblleft{}to pick up (from the ground)\textquotedblright{}.

\begin{grammarlens}[title={What you've built}]
One more word for this chapter.
\end{grammarlens}

\subsection*{Your turn}
\begin{itemize}
\raggedright
  \item \emph{Say it:} \emph{peṟukkuka}
  \item \emph{Say it:} \emph{peṟukkuka}, once more
  \item \emph{Say it:} say \emph{peṟukkuka}
  \item \emph{Recall:} say the Malayalam for to drill, then the Malayalam for to leak, then say \emph{peṟukkuka} again
\end{itemize}

\begin{tcolorbox}[breakable,colback=teal!4,colframe=teal!35!black,title={Before you move on}]
What does \ml{പെറുക്കുക} mean? (\textquotedblleft{}To pick up (from the ground)\textquotedblright{}.) Say it once more.
\end{tcolorbox}
\section[\texorpdfstring{\ml{വിതയ്ക്കുക} (vitaykkuka)}{വിതയ്ക്കുക (vitaykkuka)}]{\ml{വിതയ്ക്കുക} (vitaykkuka) — to sow (seed)}

\label{lesson:ML-C273-vitaykkuka}



\begin{quote}
Before the new one: say the Malayalam for to leak, then the Malayalam for to pick up.
\end{quote}

\subsection*{You'll want to know: \ml{വിതയ്ക്കുക}}
\textbf{\ml{വിതയ്ക്കുക}} — \emph{vitaykkuka} — \textquotedblleft{}to sow (seed)\textquotedblright{}.

\begin{grammarlens}[title={What you've built}]
One more word for this chapter.
\end{grammarlens}

\subsection*{Your turn}
\begin{itemize}
\raggedright
  \item \emph{Say it:} \emph{vitaykkuka}
  \item \emph{Say it:} \emph{vitaykkuka}, once more
  \item \emph{Say it:} say \emph{vitaykkuka}
  \item \emph{Recall:} say the Malayalam for to leak, then the Malayalam for to pick up, then say \emph{vitaykkuka} again
\end{itemize}

\begin{tcolorbox}[breakable,colback=teal!4,colframe=teal!35!black,title={Before you move on}]
What does \ml{വിതയ്ക്കുക} mean? (\textquotedblleft{}To sow (seed)\textquotedblright{}.) Say it once more.
\end{tcolorbox}
\section[\texorpdfstring{\ml{കൊത്തുക} (kottuka)}{കൊത്തുക (kottuka)}]{\ml{കൊത്തുക} (kottuka) — to peck}

\label{lesson:ML-C273-kottuka}



\begin{quote}
Before the new one: say the Malayalam for to pick up, then the Malayalam for to sow.
\end{quote}

\subsection*{You'll want to know: \ml{കൊത്തുക}}
\textbf{\ml{കൊത്തുക}} — \emph{kottuka} — \textquotedblleft{}to peck\textquotedblright{}.

\begin{grammarlens}[title={What you've built}]
One more word for this chapter.
\end{grammarlens}

\subsection*{Your turn}
\begin{itemize}
\raggedright
  \item \emph{Say it:} \emph{kottuka}
  \item \emph{Say it:} \emph{kottuka}, once more
  \item \emph{Say it:} say \emph{kottuka}
  \item \emph{Recall:} say the Malayalam for to pick up, then the Malayalam for to sow, then say \emph{kottuka} again
\end{itemize}

\begin{tcolorbox}[breakable,colback=teal!4,colframe=teal!35!black,title={Before you move on}]
What does \ml{കൊത്തുക} mean? (\textquotedblleft{}To peck\textquotedblright{}.) Say it once more.
\end{tcolorbox}
\section[\texorpdfstring{\ml{കറങ്ങുക} (kaṟaṅṅuka)}{കറങ്ങുക (kaṟaṅṅuka)}]{\ml{കറങ്ങുക} (kaṟaṅṅuka) — to spin, to go round}

\label{lesson:ML-C273-karannuka}



\begin{quote}
Before the new one: say the Malayalam for to sow, then the Malayalam for to peck.
\end{quote}

\subsection*{You'll want to know: \ml{കറങ്ങുക}}
\textbf{\ml{കറങ്ങുക}} — \emph{kaṟaṅṅuka} — \textquotedblleft{}to spin, to go round\textquotedblright{}.

\begin{grammarlens}[title={What you've built}]
One more word for this chapter.
\end{grammarlens}

\subsection*{Your turn}
\begin{itemize}
\raggedright
  \item \emph{Say it:} \emph{kaṟaṅṅuka}
  \item \emph{Say it:} \emph{kaṟaṅṅuka}, once more
  \item \emph{Say it:} say \emph{kaṟaṅṅuka}
  \item \emph{Recall:} say the Malayalam for to sow, then the Malayalam for to peck, then say \emph{kaṟaṅṅuka} again
\end{itemize}

\begin{tcolorbox}[breakable,colback=teal!4,colframe=teal!35!black,title={Before you move on}]
What does \ml{കറങ്ങുക} mean? (\textquotedblleft{}To spin, to go round\textquotedblright{}.) Say it once more.
\end{tcolorbox}
\section[\texorpdfstring{\ml{നുള്ളുക} (nuḷḷuka)}{നുള്ളുക (nuḷḷuka)}]{\ml{നുള്ളുക} (nuḷḷuka) — to pinch}

\label{lesson:ML-C273-nulluka}



\begin{quote}
Before the new one: say the Malayalam for to peck, then the Malayalam for to spin.
\end{quote}

\subsection*{You'll want to know: \ml{നുള്ളുക}}
\textbf{\ml{നുള്ളുക}} — \emph{nuḷḷuka} — \textquotedblleft{}to pinch\textquotedblright{}.

\begin{grammarlens}[title={What you've built}]
One more word for this chapter.
\end{grammarlens}

\subsection*{Your turn}
\begin{itemize}
\raggedright
  \item \emph{Say it:} \emph{nuḷḷuka}
  \item \emph{Say it:} \emph{nuḷḷuka}, once more
  \item \emph{Say it:} say \emph{nuḷḷuka}
  \item \emph{Recall:} say the Malayalam for to peck, then the Malayalam for to spin, then say \emph{nuḷḷuka} again
\end{itemize}

\begin{tcolorbox}[breakable,colback=teal!4,colframe=teal!35!black,title={Before you move on}]
What does \ml{നുള്ളുക} mean? (\textquotedblleft{}To pinch\textquotedblright{}.) Say it once more.
\end{tcolorbox}
